\documentclass{article}

\usepackage{../handouts}

\title{CSCI 535 Handout 06: Distances and Stability}
\date{17 November 2026}
\author{Name: \rule{3in}{0.15mm}}

\begin{document}
\maketitle

\section*{Warm-up: Some Definitions}

As a class, do this warm-up set together.  It is two parts:

Part 1: read the following definitions. This should be timed, 5 minutes.

For $p \in \Z_+$, the $L_p$ norm is defined as $||x||_p:=\sqrt[p]{x_1^p+x_2^p}$.
The associated distance function between points in the plane is $d_p \colon \R^2
\times \R^2 \to \R$ where $d_p(x,y)= ||x-y||_2 =
\sqrt{(x_1-y_1)^2+(x_2-y_2)^2}$.  When $p=1$, we call this the Manhattan
distance; when $p=2$, we call this the Euclidean distance.
We define $d_{\infty} \colon \R^2 \times \R^2 \to \R$~by
$$
    d_{\infty}(x,y) :=\lim_{p
    \to \infty} d_p(x,y)=\max_i\{ |x_i-y_i|\}.
$$
(Note that these distances can all be defined in $\R^d$ as well; what is under
the $p$-th root is a sum over all coordinates).

Given two sets, $P$ and $Q$, a \emph{relation} between~$P$ and $Q$ is a subset of $P \times Q$,
a \emph{matching} is a relation between~$P$ and $Q$ such
that each element in $P$ and each element in $Q$ appear at most once.
a \emph{correspondence} between~$P$ and $Q$ is a relation such that each element
of $P$ and each element of $Q$ appears at least once.
Finally, $M$ is
called a \emph{perfect matching} if $M$ is a matching and a correspondence.

As a mathematical object, \emph{a persistence diagram} is a finite set of points in
$\Rbar^2$, where $\Rbar=\R\cup \{-\infty, \infty\}$.
For defining distances, that is all we need to know about them!

Part 2: answer these questions as a class, together at the board (not silently
on your own papers). This should
take no more than 10 minutes.

\begin{enumerate}
    \item  Draw the unit sphere with $L_1$, $L_2$, and $L_{\infty}$ distances.
        (That is, draw the set of all points unit distance from the origin under
        each of these metrics).
    \item Provide an example of:
        \begin{itemize}
            \item a correspndence that is not a matching
            \item a matching that is not a correspondence
            \item a perfect matching
        \end{itemize}
    \item HW: prove that the topologies generated by open distance balls with
        the $d_p$ metric (pick a
        point and a distance $d$, the ball is all points $< d$ away from that
        point).
\end{enumerate}

\pagebreak
\section*{Distances Between Diagrams}

In what follows, let $P$ and $Q$ be point clouds in $\R^2$.
We treat $P \times Q$ as a bipartite graph, with
vertices~$P \sqcup Q$ and each $(p,q) \in P \times Q$ representing the edge
between $p$ and $q$.

\begin{enumerate}[itemsep=2in]
    \item Suppose we have a perfect matching between $P$ and $Q$. We say
        that the \emph{cost} of a matching  $M$ is the length of the longest edge
        in it, using the $d_{\infty}$ distance to calculate edge lengths.
        Describe would happen to that cost if we moved one point (in either point
        cloud) by a small amount (say, $\epsilon$).
    \item Define the distance between $P$ and $Q$ as the minimum cost matching
        among all perfect matchings between $P$ and $Q$. How many perfect
        matchings exist between $P$ and $Q$?
    \item Pick a $\delta >0$. Given the point cloud $P$, describe the set of all
        point clouds for which the distance to $P$ is $\leq
        \delta$.
    \newpage
    \item Given a threshold $\tau$, let $G_{\tau} \subset P \times Q$ such that
        all edges have length at most $\tau$. If $G_{\tau}$ has a
        perfect matching, then what do we know about the distance between $P$
        and $Q$?
    \item We can determine if a perfect matching in $G_{\tau}$ exists through
        the Ford--Fulkerson or Hopcroft--Karp algorithms.  These algorithms
        transform the bipartite graph matching problem into a flow problem.
        Describe how to transform a bipartite graph into a flow network. (Yes,
        you can lookup FF and HK if you have not seen them before, or forget
        what they are).
    \item You may have heard that the `diagonal' in the persistence diagram is
        special, and, indeed it is.
        A point $(x,x)$ would represent a topological feature that is born and
        dies at $x$: so, it exists for a trivial interval of time. Small intervals
        $(x,x+\epsilon)$ are often noise, or at lest indistinguishable from
        noise. As such,
        colloquially, we can think of allowing
        points to match to the diagonal. What is the $d_{\infty}$ distance from
        a point to the diagonal?
    \newpage
    \item Given a matching $M$, let $M^c$ denote the set of points in $P \sqcup
        Q$ that do not appear in $M$.  Then, the $p$-th Wasserstein distance
        between persistence diagrams $P$ and $Q$ is
        $$
            W_p(P,Q):=\min_M \left( \sum_{(a,b)\in M} ||a-b||^p_{\infty}
            + \frac{1}{2^{p-1}} \sum_{a \in M^c} |a_x-a_y|^p
            \right)^{1/p}
        $$
        where $M$ ranges over all matchings between $P$ and $Q$. Defining
        $W_{\infty}(P,Q):=\lim_{p \to \infty}(P,Q)$, give a formula for
        $W_{\infty}(P,Q)$.
    \item Describe what a `small change' in a persistence diagram would look
        like under $W_{\infty}(P,Q)$.
\end{enumerate}

\pagebreak
\section*{Stability}

In the last 10-ish minutes of class, discuss the following:

\begin{theorem}[Stability Theorem]
    Let $X$ be a triangulable space and $f, g: X \to \mathbb{R}$ be two tame
    continuous functions. Let $\mathcal{D}(f)$ and $\mathcal{D}(g)$ denote the
    persistence diagrams of the sublevel set filtrations of $f$ and $g$,
    respectively. Then:
    $$
        d_B(\mathcal{D}(f), \mathcal{D}(g)) \leq \|f - g\|_\infty,
    $$
    where $\|f - g\|_\infty = \sup_{x \in
    X} |f(x) - g(x)|$.
\end{theorem}

\begin{enumerate}[itemsep=2in]
    \item Why do we want this sort of stability? Explain how this is helpful.
    \item One example of how to get persistence diagrams is to take a point
        cloud and to compute the \v{C}ech complex. The Nerve Theorem tells us
        that this is the same, up to homotopy type, as studying the union of
        balls centered at our data points.  So, the \v{C}ech filtration captures
        `the topology of' the distance function to our point set.  If I add a
        point to a given point set, how much can $|f-g|_{\infty}$ change by?
\end{enumerate}

\end{document}
